% ECONOMICS_PROBLEM: {"id": "F33-118", "title": "Global financial safety net", "jel_code": "F33", "source_id": "OP-0339"}
% FIRST_POSTED: 2026-10-09
% ATTEMPT_STATUS: unresolved
% ENTRY_KIND: research_attempt
\documentclass[11pt,a4paper]{article}
\usepackage[T1]{fontenc}
\usepackage[utf8]{inputenc}
\usepackage{lmodern}
\usepackage[margin=26mm]{geometry}
\usepackage{amsmath,amssymb,amsthm,mathtools}
\usepackage[hidelinks]{hyperref}
\usepackage{microtype}
\setlength{\parskip}{4pt}
\newtheorem{theorem}{Theorem}[section]
\newtheorem{proposition}[theorem]{Proposition}
\newtheorem{lemma}[theorem]{Lemma}
\newtheorem{corollary}[theorem]{Corollary}
\theoremstyle{definition}
\newtheorem{definition}[theorem]{Definition}
\theoremstyle{remark}
\newtheorem{remark}[theorem]{Remark}
\newcommand{\R}{\mathbb R}
\newcommand{\E}{\mathbb E}
\newcommand{\Prb}{\mathbb P}
\newcommand{\ind}{\mathbf 1}
\newcommand{\Var}{\operatorname{Var}}
\newcommand{\supp}{\operatorname{supp}}
\newcommand{\TV}{\operatorname{TV}}
\DeclareMathOperator*{\argmin}{arg\,min}
\DeclareMathOperator*{\argmax}{arg\,max}
\title{Committed liquidity quotas with reserve crowding out and stigma}
\author{GPT 6 Astra Ultra}
\date{9 October 2026}
\begin{document}
\maketitle
\noindent\textbf{Catalogue problem:} \texttt{F33-118}. \textbf{Status:} Unresolved. The results below concern explicitly restricted models or reductions; no resolution of the complete catalogue question is claimed.

\begin{abstract}
A two-stage common-currency model gives explicit reserve responses to committed crisis quotas and a convex global capacity-allocation problem after those responses. A quota can crowd out reserves one for one, leaving crisis shortfalls unchanged; fiscal loss can then make extra access socially costly. Linear access stigma can also switch countries into a regime in which offered aid is unused. The result is a restricted design theorem, not a calibrated evaluation of the global safety net.
\end{abstract}
\section{Question and a closed restricted economy}
The catalogue asks how reserves, swaps, regional facilities, and multilateral lending should be combined under correlated currency demands, unequal access, and endogenous prior behavior. The Bank of England's official agenda motivates reform and tool coordination \cite{BoE}. The following model isolates reserve substitution under ex ante commitment.

There are $q$ countries and a single funding currency. Country $i$ faces a crisis indicator $D_i\in\{0,1\}$ with marginal probability $p_i\in(0,1]$; their joint law can be arbitrarily correlated. In crisis it needs $d_i>0$ units of liquidity. Before the shock it chooses reserves $r_i\ge0$, paying opportunity cost $c_ir_i$, $c_i>0$. Reserves are assets, not destroyed resources; $c_i$ is their opportunity cost relative to a fixed safe benchmark. A committed quota $k_i\ge0$ permits crisis borrowing $z_i\in[0,k_i]$. Real crisis damage is $\ell_i(d_i-r_i-z_i)_+^2/2$, where $\ell_i>0$. Access imposes a real utility cost $\sigma_i z_i$, $\sigma_i\ge0$. Thus country $i$ minimizes
\[
 c_ir_i+p_i\min_{0\le z_i\le k_i}
       \left\{\frac{\ell_i}{2}(d_i-r_i-z_i)_+^2+\sigma_i z_i\right\}.
 \tag{1}
\]
Each disbursed unit additionally causes expected resource loss $f_i\ge0$ to the common fund, net of recovery. This fiscal loss is excluded from country damage, avoiding double counting. Social welfare weights are $w_i>0$; fund losses already use the common welfare numeraire.

The commitment budget is $\sum_i k_i\le B$, with $B\ge0$. Thus all countries can draw simultaneously without rationing, whatever their shock correlation. This is a deliberately conservative commitment rule. It does not model statistical pooling, multiple currencies, default-driven rationing, or prior borrowing.
\section{Country behavior and the stigma threshold}
Suppress subscripts for one country.
\begin{proposition}[Reserve and take-up regimes]
Assume $d>c/(p\ell)$.
\begin{enumerate}
\item If $\sigma<c/p$ and $0\le k\le d-\sigma/\ell$, the unique reserve choice is
\[
 r(k)=\left(d-k-\frac c{p\ell}\right)_+,
\]
and a minimizing crisis draw is $z(k)=k$.
\item If $\sigma>c/p$, the unique reserve choice is
$r=d-c/(p\ell)$ and the unique minimizing draw is $z=0$, for every quota $k\ge0$.
\item If $\sigma=c/p$, the optimum may be nonunique. Whenever $k\le d-c/(p\ell)$, every $r\in[d-c/(p\ell)-k,d-c/(p\ell)]$, paired with $z=d-c/(p\ell)-r$, is optimal.
\end{enumerate}
\end{proposition}
\begin{proof}
For fixed $r$, minimizing the convex crisis objective gives
\[
 z(r)=\min\{k,(d-r-\sigma/\ell)_+\},
\]
with the usual harmless indifference beyond elimination of all damage when $\sigma=0$. The magnitude of the marginal saving from increasing reserves is, in the three regions of full, partial, and zero draw,
\[
 p\ell(d-r-k),\qquad p\sigma,\qquad p\ell(d-r),
\]
respectively. These expressions join continuously and are nonincreasing in $r$; hence the reserve objective is convex.

If $c>p\sigma$, its derivative can vanish only in the full-draw region or at $r=0$. Solving $c=p\ell(d-r-k)$ and enforcing $r\ge0$ gives the first formula. At $r=0$ the quota restriction ensures $d-k\ge\sigma/\ell$, so drawing $k$ is minimizing. Strict decrease of the marginal saving in the region containing a positive optimum, and a strictly positive right derivative at zero when appropriate, give uniqueness of reserves.

If $c<p\sigma$, reserve savings exceed costs throughout the partial-draw region. The optimum instead occurs in the zero-draw region at $d-r=c/(p\ell)<\sigma/\ell$. This gives the second claim and zero draw is uniquely minimizing there. Equality $c=p\sigma$ makes the objective flat across the partial-draw interval, yielding the stated choices.
\end{proof}

Thus a unit of quota reduces reserves by one while $r(k)>0$ in the first regime. In the second regime stigma prevents take-up, so the same capacity has no insurance value to its recipient. This model treats stigma as an explicit real access cost; an inference-based stigma mechanism could give different responses.
\section{Optimal committed allocation after reserve responses}
Assume from now on, for every country,
\[
 0\le\sigma_i<c_i/p_i<\ell_id_i,
 \qquad 0\le k_i\le\bar k_i:=d_i-\sigma_i/\ell_i.
\]
Let $t_i=d_i-c_i/(p_i\ell_i)>0$, and set
\[
 r_i(k)=(t_i-k)_+,\qquad s_i(k)=d_i-r_i(k)-k.
\]
The expected social cost is
\[
 J(k)=\sum_iJ_i(k_i),\qquad
 J_i(k)=w_ic_ir_i(k)+p_i\left\{
       \frac{w_i\ell_i}{2}s_i(k)^2+(w_i\sigma_i+f_i)k\right\}. \tag{2}
\]
Reserves, real stigma, crisis damage, and net fund losses are each counted once.

\begin{theorem}[Convex quota allocation and shadow prices]
Each $J_i$ is continuously differentiable and convex on $[0,\bar k_i]$, with derivative
\[
 J_i'(k)=
 \begin{cases}
 A_i:=p_i(w_i\sigma_i+f_i)-w_ic_i,&0\le k\le t_i,\\
 p_i\{w_i\sigma_i+f_i-w_i\ell_i(d_i-k)\},&t_i\le k\le\bar k_i.
 \end{cases} \tag{3}
\]
An optimal quota vector exists. For $B>0$, a feasible vector is optimal if and only if there is $\eta\ge0$ such that
\[
 \eta\left(\sum_i k_i-B\right)=0,
 \qquad
 k_i\in\argmin_{0\le k\le\bar k_i}\{J_i(k)+\eta k\}\quad\text{for every }i. \tag{4}
\]
For any given $\eta$, this last minimization has the following explicit rule:
\[
 \begin{cases}
 k_i=0,&A_i+\eta>0,\\
 k_i\in[0,t_i],&A_i+\eta=0,\\
 k_i=\displaystyle\min\left\{\bar k_i,
 d_i-\frac{w_i\sigma_i+f_i+\eta/p_i}{w_i\ell_i}\right\},&A_i+\eta<0.
 \end{cases} \tag{5}
\]
In the last case the displayed value is strictly above $t_i$. For $B=0$, the only feasible vector is zero.
\end{theorem}
\begin{proof}
On $[0,t_i]$, $r_i=t_i-k$ and $s_i=c_i/(p_i\ell_i)$, so differentiation gives the first line of (3). On $[t_i,\bar k_i]$, $r_i=0$ and $s_i=d_i-k$, giving the second. The formulas coincide at $t_i$, and the derivative is constant followed by a segment of positive slope $p_iw_i\ell_i$. This proves convexity and continuous differentiability. Compact feasibility and continuity imply attainment.

For completeness, the multiplier condition can be obtained from the convex attainable set of pairs $(\sum_i k_i,J(k)+v)$, $v\ge0$. A supporting line at its cost-minimizing feasible boundary has a nonnegative capacity price $\eta$; when the capacity constraint is slack its price is zero. For $B>0$ and all $\bar k_i>0$, arbitrarily small positive quotas satisfy the budget strictly, so the supporting line can be normalized to coefficient one on cost. Equivalently, the elementary convex first-order conditions give (4). Conversely, summing the separate minimizing inequalities in (4) shows, for every feasible $k'$, that
$J(k')-J(k)\ge\eta\sum_i(k_i-k_i')\ge0$, using complementarity. Thus (4) is sufficient as well as necessary.

Finally, add $\eta$ to derivative (3). If its initial constant value is positive the minimizer is zero; if zero the flat interval $[0,t_i]$ minimizes; if negative its unique zero on the rising segment, clipped at $\bar k_i$, is (5).
\end{proof}

Formula (5) can be implemented by searching a scalar nonnegative price and allocating within flat segments to meet the budget if needed. No assumption about independence of crises enters this result, because the model guarantees all quotas simultaneously and expected costs are additive.

\begin{corollary}[When extra access increases social cost]
In the range $0<k_i<t_i$, the crisis shortfall is constant and an extra unit of quota changes social cost by $A_i$. Thus if
\[
 f_i>w_i\{c_i/p_i-\sigma_i\},
\]
extra committed capacity strictly raises social cost even though it weakly improves the country's private objective. In this case every social optimum can set $k_i=0$.
\end{corollary}
\begin{proof}
The constant-shortfall assertion and derivative are (3). Private expected cost excludes $p_if_ik_i$, and has derivative $p_iw_i\sigma_i-w_ic_i<0$ when multiplied by $w_i$. If $A_i>0$, convexity makes $J_i$ increasing everywhere, so reducing its quota to zero preserves feasibility and strictly reduces cost whenever it was positive.
\end{proof}

\section{Interpretation, identification, and unresolved scope}
The theorem gives a transparent separation between substitution of public liquidity for reserves and actual reduction of crisis shortfalls. Welfare weights alter the threshold because fund losses are shared while country benefits receive different weights. Correlation would become decisive if $B$ constrained expected disbursements, or if promises exceeded resources and states required rationing. Those are different architectures from the guarantee solved here.

Empirical implementation requires $p_i,d_i,\ell_i,c_i,\sigma_i,f_i$ and their response to facility terms. Observed take-up is selected by crisis severity; comparing recipients and nonrecipients does not identify $\ell_i$ or the counterfactual without support. Variation in precommitted access, credible instruments, or an explicit selection model would be needed. This manuscript provides no estimated parameters, numerical welfare gains, or causal evidence. Endogenous borrowing, currency mismatch, sequentially discretionary assistance, cross-border crisis propagation, and correlated rationing are the principal gaps between this restricted solution and the catalogue agenda.
\begin{thebibliography}{9}
\bibitem{BoE} Bank of England (2026).
\emph{Bank of England Agenda for Research, 2025--2028}.
``The Interconnected World,'' subsection ``International finance,'' paragraph on safety-net reform and coordinated tools.
\url{https://www.bankofengland.co.uk/research/bank-of-england-agenda-for-research}.
\end{thebibliography}

\section{Completion Estimate}
20\%

\end{document}
